% =============================================================================
% Chapter: ER Model & Database Design
% =============================================================================

\chapter{ER Model \& Database Design}

\begin{formulabox}[Key Concepts - MEMORIZE!]
\textbf{ER Diagram Components:}
\begin{itemize}
    \item \textbf{Entity}: Rectangle - αντικείμενο/οντότητα (π.χ. Student, Course)
    \item \textbf{Attribute}: Oval - ιδιότητα οντότητας (π.χ. name, id)
    \item \textbf{Relationship}: Diamond - σύνδεση μεταξύ entities (π.χ. enrolls)
    \item \textbf{Primary Key}: Υπογραμμισμένο attribute
\end{itemize}

\textbf{Cardinality Notation:}
\begin{center}
\begin{tabular}{ll}
1:1 & One-to-One (π.χ. Person -- Passport) \\
1:N & One-to-Many (π.χ. Department -- Employee) \\
M:N & Many-to-Many (π.χ. Student -- Course) \\
\end{tabular}
\end{center}
\end{formulabox}

% =============================================================================
\section{Βασικές Έννοιες}

\subsection{Entity Types}

\begin{definitionbox}[Strong vs Weak Entity]
\begin{center}
\begin{tikzpicture}
    % Strong Entity
    \node[draw=success, fill=success!20, minimum width=2.5cm, minimum height=1cm] (strong) at (0,0) {Employee};
    \node[anchor=north, font=\small] at (strong.south) {Strong Entity};
    \node[anchor=south, font=\footnotesize\color{success}] at ([yshift=0.2cm]strong.north) {Έχει δικό του PK};

    % Weak Entity
    \node[draw=accent, fill=accent!20, minimum width=2.5cm, minimum height=1cm, double] (weak) at (6,0) {Dependent};
    \node[anchor=north, font=\small] at (weak.south) {Weak Entity};
    \node[anchor=south, font=\footnotesize\color{accent}] at ([yshift=0.2cm]weak.north) {Εξαρτάται από Strong};
\end{tikzpicture}
\end{center}
\end{definitionbox}

\subsection{Attribute Types}

\begin{center}
\begin{tikzpicture}
    % Simple
    \node[draw=secondary, ellipse, minimum width=2cm] (simple) at (0,2) {name};
    \node[anchor=north, font=\footnotesize] at (simple.south) {Simple};

    % Composite
    \node[draw=secondary, ellipse, minimum width=2cm] (comp) at (4,2) {address};
    \node[draw=secondary, ellipse, minimum width=1.5cm, font=\scriptsize] (street) at (3,0.5) {street};
    \node[draw=secondary, ellipse, minimum width=1.5cm, font=\scriptsize] (city) at (5,0.5) {city};
    \draw (comp) -- (street);
    \draw (comp) -- (city);
    \node[anchor=north, font=\footnotesize] at (4,0) {Composite};

    % Multivalued
    \node[draw=secondary, ellipse, double, minimum width=2cm] (multi) at (8,2) {phone};
    \node[anchor=north, font=\footnotesize] at (multi.south) {Multivalued};

    % Derived
    \node[draw=secondary, ellipse, dashed, minimum width=2cm] (derived) at (12,2) {age};
    \node[anchor=north, font=\footnotesize] at (derived.south) {Derived};
\end{tikzpicture}
\end{center}

% =============================================================================
\section{Cardinality \& Participation}

\begin{formulabox}[Cardinality Constraints]
\begin{center}
\begin{tikzpicture}[node distance=3cm]
    % 1:1
    \node[draw=primary, minimum width=2cm, minimum height=0.8cm] (e1) at (0,0) {Person};
    \node[draw=primary, diamond, aspect=2] (r1) at (3,0) {has};
    \node[draw=primary, minimum width=2cm, minimum height=0.8cm] (e2) at (6,0) {Passport};
    \draw (e1) -- node[above] {1} (r1);
    \draw (r1) -- node[above] {1} (e2);
    \node[anchor=west, font=\footnotesize\color{secondary}] at (7,0) {1:1 - One to One};

    % 1:N
    \node[draw=primary, minimum width=2cm, minimum height=0.8cm] (d1) at (0,-1.5) {Department};
    \node[draw=primary, diamond, aspect=2] (r2) at (3,-1.5) {employs};
    \node[draw=primary, minimum width=2cm, minimum height=0.8cm] (d2) at (6,-1.5) {Employee};
    \draw (d1) -- node[above] {1} (r2);
    \draw (r2) -- node[above] {N} (d2);
    \node[anchor=west, font=\footnotesize\color{secondary}] at (7,-1.5) {1:N - One to Many};

    % M:N
    \node[draw=primary, minimum width=2cm, minimum height=0.8cm] (s1) at (0,-3) {Student};
    \node[draw=primary, diamond, aspect=2] (r3) at (3,-3) {enrolls};
    \node[draw=primary, minimum width=2cm, minimum height=0.8cm] (s2) at (6,-3) {Course};
    \draw (s1) -- node[above] {M} (r3);
    \draw (r3) -- node[above] {N} (s2);
    \node[anchor=west, font=\footnotesize\color{secondary}] at (7,-3) {M:N - Many to Many};
\end{tikzpicture}
\end{center}
\end{formulabox}

\subsection{Participation Constraints}

\begin{center}
\begin{tikzpicture}
    % Total Participation
    \node[draw=primary, minimum width=2cm, minimum height=0.8cm] (emp) at (0,0) {Employee};
    \node[draw=primary, diamond, aspect=2] (works) at (4,0) {works\_for};
    \node[draw=primary, minimum width=2cm, minimum height=0.8cm] (dept) at (8,0) {Department};
    \draw[double] (emp) -- (works);
    \draw (works) -- (dept);

    \node[anchor=north, align=center, font=\small] at (2,-0.8) {\textbf{Total} (double line)\\Κάθε employee ΠΡΕΠΕΙ\\να δουλεύει σε dept};
    \node[anchor=north, align=center, font=\small] at (6,-0.8) {\textbf{Partial} (single line)\\Dept ΜΠΟΡΕΙ να μην\\έχει employees};
\end{tikzpicture}
\end{center}

% =============================================================================
\section{Solved Exam Problems}

\subsection{2024-25 Exam Style: ER to Relational}

\begin{examplebox}[Μετατροπή ER σε Relational Schema]
\textbf{Πρόβλημα}: Δίνεται το παρακάτω ER diagram. Μετατρέψτε σε relational schema.

\begin{center}
\begin{tikzpicture}[node distance=2.5cm]
    % Entities
    \node[draw=primary, minimum width=2cm, minimum height=0.8cm] (student) {Student};
    \node[draw=primary, diamond, aspect=2, right=of student] (enrolls) {enrolls};
    \node[draw=primary, minimum width=2cm, minimum height=0.8cm, right=of enrolls] (course) {Course};

    % Attributes for Student
    \node[draw=secondary, ellipse, above=0.8cm of student, font=\small] (sid) {\underline{sid}};
    \node[draw=secondary, ellipse, left=0.5cm of student, font=\small] (sname) {name};

    % Attributes for Course
    \node[draw=secondary, ellipse, above=0.8cm of course, font=\small] (cid) {\underline{cid}};
    \node[draw=secondary, ellipse, right=0.5cm of course, font=\small] (cname) {title};

    % Relationship attribute
    \node[draw=secondary, ellipse, below=0.5cm of enrolls, font=\small] (grade) {grade};

    % Connections
    \draw (student) -- node[above] {M} (enrolls);
    \draw (enrolls) -- node[above] {N} (course);
    \draw (sid) -- (student);
    \draw (sname) -- (student);
    \draw (cid) -- (course);
    \draw (cname) -- (course);
    \draw (grade) -- (enrolls);
\end{tikzpicture}
\end{center}

\textbf{Λύση - Βήμα προς βήμα:}

\textbf{Step 1: Strong Entities → Tables}
\begin{itemize}
    \item \texttt{Student(\underline{sid}, name)}
    \item \texttt{Course(\underline{cid}, title)}
\end{itemize}

\textbf{Step 2: M:N Relationship → New Table}
\begin{itemize}
    \item \texttt{Enrolls(\underline{sid}, \underline{cid}, grade)}
    \item Foreign Keys: sid → Student, cid → Course
\end{itemize}

\textbf{Final Schema:}
\begin{center}
\begin{tikzpicture}
    \node[draw=success, fill=success!20, rounded corners=5pt, minimum width=10cm, minimum height=3cm] (schema) {};
    \node[anchor=center, align=left, font=\ttfamily\small] at (schema.center) {
        Student(\underline{sid}, name)\\
        Course(\underline{cid}, title)\\
        Enrolls(\underline{sid}, \underline{cid}, grade)\\
        \hspace{1cm}FK: sid REFERENCES Student(sid)\\
        \hspace{1cm}FK: cid REFERENCES Course(cid)
    };
\end{tikzpicture}
\end{center}
\end{examplebox}

\subsection{Cardinality Rules for Conversion}

\begin{center}
\begin{tabular}{lll}
\toprule
\textbf{Cardinality} & \textbf{Conversion Rule} & \textbf{Example} \\
\midrule
1:1 & FK σε οποιαδήποτε πλευρά & Person-Passport \\
1:N & FK στη N-πλευρά & Dept-Employee \\
M:N & Νέος πίνακας με 2 FKs & Student-Course \\
\bottomrule
\end{tabular}
\end{center}

% =============================================================================
\section{Common Mistakes}

\begin{warningbox}[Συχνά Λάθη στις Εξετάσεις]
\begin{enumerate}
    \item \textbf{M:N χωρίς νέο πίνακα}:
    \begin{itemize}
        \item[$\times$] Βάζω FK στη μία πλευρά
        \item[\checkmark] Δημιουργώ junction table με composite key
    \end{itemize}

    \item \textbf{Weak Entity conversion}:
    \begin{itemize}
        \item[$\times$] Μόνο το partial key ως PK
        \item[\checkmark] Composite key: (owner\_PK + partial\_key)
    \end{itemize}

    \item \textbf{Multivalued attributes}:
    \begin{itemize}
        \item[$\times$] Πολλές στήλες (phone1, phone2, ...)
        \item[\checkmark] Ξεχωριστός πίνακας με FK
    \end{itemize}

    \item \textbf{Total Participation}:
    \begin{itemize}
        \item[$\times$] Αγνοώ τη διπλή γραμμή
        \item[\checkmark] NOT NULL constraint στο FK
    \end{itemize}
\end{enumerate}
\end{warningbox}

% =============================================================================
\section{Quick Reference}

\begin{center}
\begin{tikzpicture}
    \node[draw=ntua-blue, fill=ntua-blue!10, rounded corners=5pt, minimum width=14cm, minimum height=6cm] (main) {};

    \node[anchor=north west, font=\bfseries\Large\color{ntua-blue}] at ([xshift=5pt, yshift=-5pt]main.north west) {ER → Relational Conversion Rules};

    \node[anchor=north west, align=left, font=\small] at ([xshift=10pt, yshift=-35pt]main.north west) {
        \textbf{1. Strong Entity} → Table με όλα τα simple attributes\\[3pt]
        \textbf{2. Weak Entity} → Table με (owner\_PK + partial\_key) ως composite PK\\[3pt]
        \textbf{3. 1:1 Relationship} → FK σε οποιαδήποτε πλευρά (προτίμηση: total participation)\\[3pt]
        \textbf{4. 1:N Relationship} → FK στην N-πλευρά (many side)\\[3pt]
        \textbf{5. M:N Relationship} → Νέος πίνακας με PKs και των δύο entities\\[3pt]
        \textbf{6. Multivalued Attr} → Ξεχωριστός πίνακας (entity\_PK, value)\\[3pt]
        \textbf{7. Composite Attr} → Flatten: κάθε component γίνεται στήλη\\[3pt]
        \textbf{8. Derived Attr} → ΔΕΝ αποθηκεύεται (υπολογίζεται)
    };
\end{tikzpicture}
\end{center}
